% ---------------------------------------------------------------------------
% 2 Background
% Source: docs/shared/backend_selection.md, docs/overview.md,
%         docs/hybrid/t123_walkthrough.md §3 (glossary) and §3.1
% ---------------------------------------------------------------------------

\section{Background}

\subsection{DNAGPT and the genomic task families}
DNAGPT is an autoregressive transformer adapted to sequence and numerical tasks over
DNA~\cite{dnagpt2023}. We evaluate its released $0.1$-billion-parameter configuration. The backbone has
$12$ transformer blocks, hidden width $768$, and $12$ causal-attention heads of width $64$. Its
feed-forward sublayer expands to width $3{,}072$. DNA is represented through $6$-mer tokenization
and combined with task-format tokens before entering the transformer.

Each block follows the structure that determines the encrypted circuit: LayerNorm; query, key, and
value projections; causal self-attention; an output projection and residual; a second LayerNorm;
then an expanded MLP with GELU and another residual. The dense maps apply fixed model weights.
LayerNorm and GELU apply nonlinear functions to the activations, while attention introduces products
between two data-dependent quantities and normalizes them with a softmax.

The release includes fine-tuned heads for genomic-signal recognition and mRNA abundance regression,
but none for the genome-understanding benchmark, so Section~\ref{sec:baseline} evaluates that
backbone with a locally fine-tuned linear head. All three task families predate DNAGPT and carry
established task-specific baselines~\cite{deepgsr,xpresso,dnabert2}, against which
Section~\ref{sec:baseline} establishes model fidelity; the encrypted study then uses the
signal-recognition path as its task-length anchor.

\subsection{CKKS}
CKKS supports approximate arithmetic on packed vectors of real or complex
values~\cite{ckks2017}. Encoding places many values into the slots of one plaintext, which is then
encrypted under the data owner's key. Addition and multiplication act slotwise across the packed
vector, while rotations move values between slots. This single-instruction, multiple-data model is
a natural match for transformer activations: a rotation-and-diagonal schedule can express a public
matrix transform as ciphertext--plaintext multiplication and accumulation.

Not every transformer product has a public operand. Query--key scoring multiplies two encrypted
projections of the input, and attention context multiplies encrypted attention weights by encrypted
values. Those operations require ciphertext--ciphertext multiplication and relinearization. The
distinction is about data ownership rather than implementation convenience: public weights and
masks remain plaintext, while every query-derived operand remains encrypted at the server.

CKKS controls approximate fixed-point scale through rescaling. A multiplication consumes a level
from a finite modulus chain, and the computation fails once that multiplicative budget is exhausted.
Bootstrapping can refresh a ciphertext without decryption, while fresh client encryption can reset
the budget in an interactive protocol. Both choices carry system costs, so multiplicative depth is
an architectural resource rather than only a cryptographic parameter.

\subsection{Why the nonlinearities are the architecture}
CKKS directly supplies polynomial arithmetic, but a transformer requires inverse square root in
LayerNorm, exponentiation and normalization in softmax, and GELU. A non-interactive circuit must
replace these functions with polynomial approximations over declared input ranges. Higher-degree
approximations consume additional multiplicative depth; across repeated blocks, the evaluator must
provide a deeper context, a bootstrapping schedule, or both. Larger modulus chains and refresh
operations also increase the context and evaluation-key material that must remain resident.

The alternative is to place exact functions at declared client boundaries, which removes
approximation-range calibration and resets the encrypted depth budget but requires an online client
and exposes intermediate activations to it. The choice between polynomial approximation and client
evaluation therefore determines memory, interaction, depth, and the threat model.
Section~\ref{sec:protocol} evaluates both designs and reports where each stops.

\subsection{Evaluated software backends}
We separate correctness semantics from the performance path. OpenFHE provides CKKS context and
security-parameter construction, encoding, rotations, polynomial evaluation, and approximate
bootstrapping~\cite{openfhe}. It supplies the reference behavior used to validate local encrypted
operators and the non-interactive design. The evaluated GPU backend interoperates with those
representations and provides CUDA implementations of the server-side CKKS operations; all reported
accelerator measurements use that path. Python, NumPy, and PyTorch construct and validate the frozen
plaintext reference, but they do not participate in encrypted performance measurements.

We also evaluated the Boolean/integer-oriented TFHE-rs path and Concrete ML's client/server
workflow~\cite{concreteml,tfhers}; their available GPU acceleration and ciphertext expansion were
less suitable for this matrix-heavy graph than the already validated CKKS path, and a separate CKKS
wrapper lacked the refresh and general nonlinear-function support the non-interactive baseline
needs. These decisions select among the software paths tested here. They do not establish that CKKS
is the only scheme able to express encrypted transformer inference.
